\changecolours{ScoutGreen}
\section[Conclusão e Próximos Passos]{Checklist do Entregável CASE e Ponte para a Aula 09}

% ------------------------------------------------------------------------------
\twocolframe[title={Checklist Final de Qualidade do DER/EER (AVP1)},leftcol=7.0cm,rightcol=7.0cm]{%
  \alert{Antes de Exportar o PDF/PNG, Verifique:}\par
  \itemseps[topsepi=0.5mm,itemsepi=1mm,labelsep=2mm,leftmargini=3mm]
  \begin{itemize}
    \item[$\boxtimes$] Todas as entidades fortes possuem chave primária sublinhada ($\underline{\text{pk}}$) ou pino preto ($\bullet$)?
    \item[$\boxtimes$] A entidade fraca possui retângulo duplo, losango duplo, linha dupla e discriminador tracejado?
    \item[$\boxtimes$] A hierarquia EER explicita as duas restrições ($d/o$ e Total/Parcial ou $t/p, e/c$)?
    \item[$\boxtimes$] O auto-relacionamento possui os papéis (\textit{roles}) escritos nas duas arestas?
    \item[$\boxtimes$] Não há nenhuma chave estrangeira desenhada como atributo no DER?
  \end{itemize}
}{%
  \alert{Ponte Pedagógica: Rumo à Aula 09!}\par
  \itemseps[topsepi=0.5mm,itemsepi=1mm,labelsep=2mm,leftmargini=3mm]
  \begin{itemize}
    \item \textbf{Fechamento do Ciclo Conceitual:} Com a Aula 08, concluímos a modelagem conceitual (Unidade III).
    \item \textbf{Aula 09 (21/10/2026) -- O Modelo Relacional (Edgar F. Codd, 1970):} Estudaremos domínios atômicos, tuplas, esquemas de relação, superchaves, chaves candidatas, chave primária (PK), chave estrangeira (FK) e as regras formais de integridade!
  \end{itemize}
}
